\section{Discussion}
\label{sec:discussion}

\paragraph{Why the tradeoff disappears on CPSC2018.}
The most striking finding is the complete disappearance of the discrimination--calibration
tradeoff on CPSC2018. We attribute this to three structural differences from PTB-XL.
First, task type: CPSC2018 is a single-label 9-class task (each record has one dominant
arrhythmia), while PTB-XL is a 5-class multi-label task with correlated superclasses.
Multi-label calibration requires the linear probe to simultaneously assign valid probability
mass to correlated outputs --- a harder geometric problem that amplifies systematic
over- or under-confidence.
Second, class overlap: PTB-XL superclasses (e.g., MI and CD) frequently co-occur,
forcing the probe to partition probability across correlated classes. CPSC2018 classes
are more mutually exclusive.
Third, class prevalence: PTB-XL has highly variable class prevalences
(NORM: 48\%, HYP: 8\%), requiring the probe to operate at different operating points
for each class simultaneously. These structural differences suggest that calibration
auditing is most important for multi-label cardiac classification tasks.

\paragraph{Pretraining objective mechanism.}
MERL's underconfidence ($\text{SCE}=-0.003$ on PTB-XL) is consistent with the
documented behaviour of contrastive vision FMs~\citep{hekler2025beyond}. The mechanistic
interpretation is that contrastive objectives (InfoNCE) push representations apart to
maximise inter-class separation, producing a norm-uniform representation space with high
feature spread (std of L2 norms: 34.8 vs 0.23 for ST-MEM). When this spread representation
is projected through a linear probe, the probability outputs may be diffuse and underconfident
--- the probe cannot confidently assign high probabilities because the representation space
does not concentrate examples in tight class clusters.

The two masked-prediction models behave similarly in \emph{direction} but differ in
\emph{magnitude} of miscalibration. ST-MEM (masked patch reconstruction) and HuBERT-ECG
(masked hidden-unit prediction) both have near-zero SCE on PTB-XL ($-0.0001$ and $-0.0007$),
indicating no systematic over/under-confidence bias; yet both have the largest \emph{unsigned}
calibration error (ECE 0.124 and 0.136). Feature geometry explains the distinction: HuBERT-ECG
has the highest cosine class separation (0.046) of any model, and ST-MEM the second highest
(0.017), whereas MERL's separation is near zero (0.004). High class separation concentrates
predictions near 0/1, so when the empirical frequency disagrees the per-bin gap is large ---
producing high ECE even when the signed bias cancels out. Masked-prediction objectives thus
yield ``confidently wrong in both directions'' probabilities that recalibration is especially
effective at fixing (HuBERT-ECG ECE drops 71\% under isotonic).

ECGFM-KED (knowledge-enhanced, $+0.007$), ECG-FM (contrastive+generative, $+0.008$), and the
supervised S4D baseline ($+0.007$) are mildly overconfident on PTB-XL, suggesting that
decoder/generative and discriminative-supervised training shift calibration toward
overconfidence. This graded pattern --- contrastive underconfident, masked-prediction
neutral-but-high-ECE, knowledge/generative/supervised overconfident --- motivates further
study of how pretraining objectives shape probability calibration.

\paragraph{Recalibration recommendation.}
Based on the calsize analysis and DCA results, we summarise deployment recommendations:

\begin{center}
\begin{tabular}{lll}
\toprule
Calibration set size & Recommended method & Rationale \\
\midrule
$n_{\mathrm{cal}} < 63$      & Platt scaling          & Isotonic overfits; Platt is stable \\
$63 \leq n_{\mathrm{cal}} < 200$ & Isotonic regression    & Better ECE; crossover confirmed at $n=63$ \\
$n_{\mathrm{cal}} \geq 200$  & Isotonic regression    & Consistent 43--64\% ECE reduction \\
\bottomrule
\end{tabular}
\end{center}

Temperature scaling is not recommended for multi-label ECG tasks: it reduces ECE by only
11--38\% on PTB-XL, consistently worse than Platt scaling. The MLP head is unreliable
at $n_{\mathrm{cal}} < 100$, where it frequently underfits.

\paragraph{MIMIC-IV-ECG and within-domain pretraining.}
The MIMIC-IV-ECG results reveal a counter-intuitive finding: ECG-FM, which was pretrained
on MIMIC-IV-ECG data, achieves AUROC 0.771 on MIMIC-IV-ECG evaluation---below MERL (0.861)
and ST-MEM (0.854) in every resplit, and statistically tied with the supervised S4D baseline
(0.793; below S4D in only 2/10 seeds). This demonstrates that
backbone pretraining domain alignment is a necessary but not sufficient condition for
downstream linear-probe performance. The phenomenon likely arises because ECG-FM's hybrid
contrastive+generative pretraining objective optimises for self-supervised tasks (masked
prediction, contrastive similarity) that do not directly align with the
classification targets derived from machine-measurement reports. By contrast, MERL's
contrastive pretraining with clinical text descriptions may produce representations
more naturally aligned with the lexical structure of the machine measurement labels.
This has a practical implication: practitioners should not assume that a model pretrained
on their target cohort (e.g., in-house MIMIC-IV deployment) will automatically be the best
choice for a given classification task without linear-probe evaluation.

\paragraph{Cross-dataset calibration transfer.}
ECE varies by a factor of up to $10.3\times$ for the same model across datasets
(ST-MEM: CODE-15\% ECE 0.012 vs PTB-XL ECE 0.124). The MIMIC-IV-ECG dataset shows
intermediate calibration (ECE 0.033--0.061), suggesting that the calibration difficulty
correlates with label structure: binary single-class tasks (MIMIC machine labels, CODE-15\%)
are easier to calibrate than correlated multi-label superclasses (PTB-XL). Models calibrated
on CODE-15\% may be substantially miscalibrated on PTB-XL. This motivates re-calibration
on local data from the target deployment population.

\paragraph{Limitations.}
Six limitations bound this study's scope and generalisability.

\textbf{Frozen linear probes.} All evaluations use logistic linear probes on frozen
FM representations. Full fine-tuning may alter calibration --- potentially improving
or worsening it --- in ways not captured here. Calibration under fine-tuning is an
important open question.

\textbf{Sample sizes.} PTB-XL and MIMIC-IV-ECG use 2{,}000-record subsets and CODE-15\%
uses ${\sim}8{,}000$ records sampled across five HDF5 parts (chosen so that every rare
arrhythmia class clears the minimum-support threshold); the full datasets are larger still.
Larger samples may shift absolute ECE values, though the qualitative PTB-XL vs.\ arrhythmia
contrast is unlikely to reverse.

\textbf{HuBERT-ECG preprocessing.} HuBERT-ECG is included via its public \texttt{safetensors}
checkpoint. Following the authors' pipeline, the 12 leads are flattened into a single 1-D
sequence and decimated before encoding; this single-channel treatment differs from the
native multi-lead handling of the other encoders and may modestly affect its representations.

\textbf{ECGFM-KED architecture.} The partial architecture reconstruction (35 missing keys)
means ECGFM-KED results are lower bounds on true performance. The partial model likely
underestimates what the full ECGFM-KED would achieve.

\textbf{DCA scope.} The real DCA covers ST-MEM only.
Extending to MERL and ECG-FM is straightforward from cached features but not included.

\textbf{Subgroup analysis.} The age/sex subgroup analysis uses a re-trained probe
on the same 2,000-record sample; absolute calibration values are not directly
comparable to the main benchmark and the analysis is exploratory.

\paragraph{The ecg-audit package.}
The \texttt{ecg-audit} package provides these recalibration methods, plus the PhysioCalib
severity-weighted variant, as a one-command API. The package is installed and tested with
unit tests covering all recalibrators. The leaderboard scaffold allows the community to
populate calibration scores for additional FMs as checkpoints become available.
